% ------------------------------------------------------------------------
% file `partie1-exercise-1-solution-body.tex'
%   in folder `xsim/'
%
%     solution of type `exercise' with id `1'
%
% generated by the `solution' environment of the
%   `xsim' package v0.21 (2022/02/12)
% from source `partie1' on 2023/06/07 on line 40
% ------------------------------------------------------------------------
    Évaluation :
    \begin{itemize}
        \item Les bonnes réponses sont cochées et aucune mauvaise réponse n'est cochée \hfill 1
    \end{itemize}
    \begin{enumerate}
        \item Une fonction linéaire est une fonction qui \dots
              \begin{QCM}
                  \choice[\faCheckSquare] passe par l'origine du repère
                  \choice a une pente nulle
                  \choice est toujours croissante
                  \choice[\faCheckSquare] admet $0$ comme racine.
              \end{QCM}
        \item Combien de solution(s) admet l'équation $x^2+4=0$ ?
              \begin{QCM}(3)
                  \choice $2$
                  \choice $1$
                  \choice[\faCheckSquare] $0$
              \end{QCM}
        \item Quel est le degré du polynôme $x^6 + 5x^4 -2x^3 +23$ ?
              \begin{QCM}(3)
                  \choice $3$
                  \choice $4$
                  \choice[\faCheckSquare] $6$
              \end{QCM}
        \item Quelle(s) expression(s) est/sont équivalentes à $\dfrac{3}{\sqrt{2}}$ ?
              \begin{QCM}(4)
                  \choice[\faCheckSquare] $\dfrac{3\sqrt{2}}{2}$
                  \choice $\dfrac{\sqrt{6}}{2}$
                  \choice $\sqrt{\dfrac{3}{2}}$
                  \choice[\faCheckSquare] $\sqrt{\dfrac{9}{2}}$
              \end{QCM}
    \end{enumerate}
